El 1971 l'astronauta David Scott, de la missió \textit{Apollo 15}, va fer l'experiment següent a la superfície de la Lluna: en una mà hi tenia una ploma de falcó de 30~g de massa i a l'altra mà hi tenia un martell d'alumini d'1,32~kg. Els va deixar anar alhora des de la mateixa altura i va comprovar la predicció de Galileu segons la qual en caiguda lliure els dos objectes havien d'arribar simultàniament a terra. Concretament, tots dos objectes van trigar 1,1~s a recórrer els 100~cm que els separaven del terra.

\begin{apartat}{1,25}
A partir de l'experiment de David Scott, calculeu la intensitat del camp gravitatori a la superfície de la Lluna i la massa de la Lluna.
\end{apartat}

\begin{apartat}{1,25}
Calculeu el període orbital de la Lluna al voltant de la Terra.
\end{apartat}

\dades{%
  $G = 6,67\times 10^{-11}\ \mathrm{N\,m^2\,kg^{-2}}$.\\
  $M_\mathrm{Terra} = 5,98\times 10^{24}\ \mathrm{kg}$.\\
  Distància Terra-Lluna $= 3,84\times 10^{8}\ \mathrm{m}$.\\
  $R_\mathrm{Lluna} = 1,74\times 10^{6}\ \mathrm{m}$.}
